\documentclass[UTF8,12pt]{ctexart}

\usepackage[a4paper,margin=2.2cm]{geometry}
\usepackage{amsmath}
\usepackage{booktabs}

\setlength{\parindent}{2em}
\setlength{\parskip}{0.36em}
\linespread{1.15}

\newcommand{\cue}[1]{%
  \par\addvspace{0.85em}%
  \noindent{\large\bfseries #1}\par\nopagebreak\smallskip\noindent%
}

\title{基于领域自适应的工业过程故障诊断方法研究\\答辩演讲稿}
\author{08022311\quad 陈鲲龙}
\date{2026年6月}

\begin{document}
\maketitle

\cue{封面页}
我是 08022311 陈鲲龙，在王永健老师的指导下完成了课题：《基于领域自适应的工业过程故障诊断方法研究》。

\cue{研究问题界定}
题目中的主语和定语都是比较大的研究方向，我这里先收敛到本毕设具体面向的问题。关于领域自适应，本文使用 TE 过程数据集，它一共包含 6 个生产工况，我把这 6 个工况看作 6 个领域；它们可以作为源域或目标域相互迁移和适应，具体组合在后面的实验方案中说明。关于故障诊断，TE 数据集中包含正常状态 0 类和 1 到 28 类故障，因此在本毕设中，故障诊断本质上被建模为一个 29 类分类问题。右上角页码对应论文页码，方便老师们对照查看。

本文要解决的核心问题是：用源工况有标签数据训练好的诊断模型，由于工况之间数据分布不同，直接用到目标工况时性能会下降；而目标工况训练阶段又没有故障标签。因此，本文研究的是无监督跨工况领域自适应故障诊断。

\cue{核心成果}
本文核心成果是 3 个创新算法，其中两个面向单源场景，一个面向多源场景，并且三者之间存在递进关系。老师们可以看到第 7 页给出了全文结构：这 3 个方法分别对应论文第三章、第四章和第五章。后面三套方法都按同一个节奏讲：架构页说明缩写对应的新增机制，公式页解释公式里各项含义，训练流程页按论文 Algorithm 说明每一步做什么。

\cue{方法一：TPU-DeepJDOT 架构}
第一套方法是 TPU-DeepJDOT。这里的 T 是 temporal，新增原始时间窗口统计代价；P 是 prototypical，新增源域类别原型和相对原型代价；U 是 unbalanced，把 DeepJDOT 原来的平衡传输改成非平衡传输。也就是说，方法一是在 DeepJDOT 的特征--标签联合传输上，额外加入时序信息、类别中心结构和非平衡匹配机制。

\cue{方法一：联合代价与非平衡传输}
这页主要解释联合代价公式。$C_{ij}$ 是源样本 $i$ 和目标样本 $j$ 的匹配代价：$c^f_{ij}$ 看深度特征距离，$c^y_{ij}$ 看目标样本被预测成源样本类别的标签代价，$r_{ij}^{\mathrm{proto}}$ 看目标样本是否相对更接近当前源类原型，$c_{ij}^{\mathrm{temp}}$ 看原始时间窗口统计量是否接近。下面的总目标 $\mathcal{L}_{\mathrm{TPU}}$ 由源域分类、源域监督对比和非平衡传输对齐三部分组成，其中非平衡传输允许削减明显不合适的匹配质量。

\cue{方法一：TPU-DeepJDOT 训练流程}
这页对应论文第三章的 Algorithm。每一轮先取一个源域有标签 batch 和一个目标域无标签 batch；FCN 提取特征并给出分类预测；然后用源域标签更新类别原型，同时计算源域分类损失和监督对比损失；再把特征、标签、原型和时序四类代价合成 $C_{ij}$，用非平衡 Sinkhorn 求传输计划 $\gamma^\ast$；最后把这些损失加权合起来更新模型。整个流程里目标标签不参与训练。

\cue{方法一验证}
从代表任务 $m_5\rightarrow m_2$ 可以看到，SourceOnly 的目标域准确率为 56.08\%，DeepJDOT 提升到 67.20\%，而 TPU-DeepJDOT 达到 83.60\%。这说明非平衡传输、类别原型和时序代价确实改善了对齐质量。t-SNE 可视化也进一步说明，方法一不仅让源域和目标域分布更接近，同时让类别结构更加清晰。

\cue{方法二：CBTPU-DeepJDOT 架构}
第二套方法是 CBTPU-DeepJDOT。C 是 confidence，表示多证据置信筛选；B 是 balanced，表示对接受的伪标签做类别均衡；T、P、U 继承方法一，仍然对应 temporal、prototypical 和 unbalanced。架构上，CBTPU 不是重做一个新主干，而是在 TPU-DeepJDOT 的基础上增加教师模型、传输语义、教师预测语义、原型语义，以及伪标签筛选和一致性训练模块。

\cue{方法二：多证据伪标签筛选}
这页解释伪标签筛选公式。$q^{\mathrm{ot}}$ 来自传输计划，表示目标样本主要接收哪些源类质量；$q^{\mathrm{cls}}$ 来自固定 TPU 教师模型；$q^{\mathrm{proto}}$ 来自目标特征到源域类别原型的距离。三者按幂加权几何平均融合成 $q^{\mathrm{mix}}$。接受条件 $\mathcal{A}_j$ 有三道门：三类证据的类别一致，或者 JS 差异足够小；融合置信度超过动态阈值；传输语义熵足够低。总目标则是在 TPU 原有损失上，增加伪标签损失、一致性损失和一个很小的信息最大化正则。

\cue{方法二：CBTPU-DeepJDOT 训练流程}
这页对应论文第四章的 Algorithm。先用第一阶段 TPU-DeepJDOT 初始化学生模型，并作为固定教师参考；每轮采样源域和目标域 batch，对目标样本生成弱增强和强增强；教师看弱增强，给出 $q^{\mathrm{cls}}$，学生同时参与源域训练、传输对齐和强增强预测；再由传输计划得到 $q^{\mathrm{ot}}$，由原型距离得到 $q^{\mathrm{proto}}$，三类证据融合成 $q^{\mathrm{mix}}$。通过筛选的目标样本进入软伪标签损失，中等置信但传输稳定的样本进入强弱增强一致性损失，最后只更新学生模型。

\cue{方法二验证}
在代表任务 $m_5\rightarrow m_1$ 上，TPU-DeepJDOT 的准确率为 77.59\%，加入伪标签细化后的 CBTPU-DeepJDOT 提升到 81.38\%。从局部 t-SNE 结果看，CBTPU-DeepJDOT 对目标域边界有进一步收紧作用。也就是说，伪标签阶段的提升幅度不像方法一那么大，但它在可靠对齐基础上继续修正了目标域决策边界。

\cue{方法三：CA-CCSR-WJDOT 架构}
第三套方法是 CA-CCSR-WJDOT。CA 表示 CoDATS-augmented，也就是加入 CoDATS 风格时间序列表征和领域对抗；CCSR 表示 class-conditional source reliability，也就是类条件源域可靠性；WJDOT 是 weighted joint distribution optimal transport，即多源加权联合分布最优传输。架构上的新增元素是源域--故障类别可靠性矩阵 $R_{kc}$ 和类别级源域权重 $\alpha_{kc}$，把 WJDOT 的源域整体权重细化到每个故障类别上。

\cue{方法三：类条件源域可靠性}
这页解释可靠性公式。$R_{kc}$ 表示第 $k$ 个源域对第 $c$ 类故障的可靠性分数，由四项加权得到：源目标原型距离 $D^{\mathrm{proto}}_{kc}$，类条件传输代价 $D^{\mathrm{ot}}_{kc}$，目标预测熵 $H^{\mathrm{pred}}_{kc}$，以及源域类别召回误差 $E^s_{kc}$。这些量越小，说明该源域在这个类别上越可靠。再对 $-R_{kc}$ 做 softmax，得到类别级源域权重 $\alpha_{kc}$。CCSR 损失就是用 $\alpha_{kc}$ 去加权各源域、各类别的传输代价；总目标还包括源域监督、领域对抗、WJDOT 聚合损失和教师先验约束。

\cue{方法三：CA-CCSR-WJDOT 训练流程}
这页对应论文第五章的 Algorithm。每轮从多个源域分别取有标签 batch，再取目标域无标签 batch；共享 FCN 和 CoDATS 分类头先提取特征并做源域分类；然后对每个源域单独和目标域求一次 WJDOT 传输，得到源域级传输损失、类条件传输质量和类条件传输代价；接着用原型距离、传输代价、预测熵和源域召回误差构造 $R_{kc}$，再转成 $\alpha_{kc}$；最后用源域监督、领域对抗、WJDOT、CCSR 和教师先验这些损失一起更新模型。训练结束后，再用教师先验保护的融合规则生成目标域预测。

\cue{方法三验证}
在五源到 $m_2$ 的代表任务中，直接合并源域的 SourceOnly 准确率为 64.02\%，CA-CCSR-WJDOT 达到 82.89\%。对应的 t-SNE 结果显示，目标域类别聚集更明显，误判结构减少。源域--类别权重矩阵也能解释这一点：不同故障类别并不依赖同一个源域，而是有的类别更依赖 $m_1$，有的更依赖 $m_4$ 或 $m_6$，也有部分类别由多个源域共同提供信息。

\cue{实验协议}
实验部分我设置了三类场景：单源到单目标任务一共 6 个，两源到单目标任务 3 个，五源到单目标任务 6 个。评价指标包括整体准确率、宏平均 F1、平衡准确率、混淆矩阵和 t-SNE 可视化。所有无监督领域自适应方法都遵守同一约束：目标域标签不参与训练、早停、模型选择或参数确定，只在最终评价时使用。

\cue{汇总结果}
从单源汇总结果看，DeepJDOT 的平均准确率为 68.88\%，TPU-DeepJDOT 达到 80.54\%，相对 DeepJDOT 提升 11.66 个百分点；CBTPU-DeepJDOT 进一步达到 82.02\%，并在 6 个单源迁移方向上都取得无监督方法第一。这说明单源性能的主提升来自更可靠的联合分布对齐，伪标签阶段则提供了稳定的边界细化。

从多源结果看，CA-CCSR-WJDOT 在 3 个两源任务上的平均准确率为 81.21\%，在 6 个五源任务上达到 85.50\%，并且两源和五源任务中都取得了无监督方法第一。相对 WJDOT，CA-CCSR-WJDOT 在多源任务上的平均准确率提升为 7.69 个百分点。配对显著性检验也表明，三组相邻方法的提升在对应迁移场景中均为正向，且 Wilcoxon $p$ 值小于 0.05。

\cue{结论页}
最后总结一下。本文围绕 TEP 跨生产模式故障诊断，形成了三条递进的方法链。TPU-DeepJDOT 解决源目标样本匹配质量问题；CBTPU-DeepJDOT 解决目标样本伪标签可靠性问题；CA-CCSR-WJDOT 解决多源场景下源域类别贡献分配问题。实验结果说明，跨工况故障诊断不能只追求整体分布对齐，而需要同时处理样本匹配、目标样本筛选和源域--类别贡献分配。以上就是我的汇报，请各位老师批评指正。

\end{document}
